\justifying
\NSFCBodyText

\subsubsection{与本项目相关的研究工作积累}
申请人颜丰为新疆大学副教授、博士、硕士研究生导师，具有多源信息表示、跨模态检索、视觉问答和软件系统研发的研究积累。与本项目直接相关的基础主要来自证据检索、条件化信息选择和长上下文效率研究。上述积累支持算法实现与实验组织，但尚不能证明已完成智能体上下文压缩或有效恢复的先导实验。

在跨模态检索方面，申请人作为通讯作者参与 Feature Fusion Mamba Hashing via Decoupling for Cross-Modal Retrieval，负责研究构思、方法设计、实验指导及论文修改\cite{yanhash}。该工作涉及信息表示与检索组织，可支持本项目比较不同证据索引和表示方式。任务模态与评价终点不同，因此项目将重新建立动作依赖及证据有效性实验，不把跨模态检索结果直接当作长程智能体性能。

在条件化信息选择和结构推理方面，申请人作为通讯作者参与 RSHR+，负责方法设计、实验指导与结果分析\cite{yanrshr}；作为第一作者兼通讯作者完成知识感知视觉问答研究\cite{yanvqa}。这些工作为研究“任务条件怎样改变信息选择”提供方法与实验经验，可迁移到本项目的状态提取、依赖关系对照与错误分析。其支持限于研究能力与相关方法，不等同于已验证本项目假说。

申请人参与“可学习乘积量化与状态空间模型协同的长上下文高效注意力机制研究”，正式负责高效注意力机制研究。该经历有助于理解长上下文资源消耗和效率评价。本申请研究运行时的上下文选择与证据恢复，以冻结执行模型开展主要比较，与改变注意力实现或模型内部结构的任务形成明确区分。

\subsubsection{已取得的研究工作成绩}
上述三项代表作的题录、作者角色与贡献已在共用申请人事实中确认，本次又核对出版元数据。它们分别支撑检索表示、条件化选择与知识增强推理；研究方案将把这些能力落实到可复现基线、成对干预及模型误用证据分析。正在承担的知识库研发工作提供数据组织和问答系统工程背景，但已有合同成果及知识产权受双方约定约束，不能未经授权移入本项目。

目前尚无可复用的本项目级数据、代码、模型配置、评价脚本或消融日志。项目启动后的首项任务是建立共同评测和基线，而不是以未经记录的实验数值支撑可行性。预期性能、成本节省和业务提升均保留为待检验目标；图件只表达研究设计、流程及条件，不展示虚构实验曲线或前期成果截图。

\subsubsection{可行性分析}
\textbf{方法可行性。}状态压缩、外部记忆、版本关联及可复位工具环境均有公开研究依据，核心挑战可以拆为可独立验证的局部问题。项目先用显式规则建立基线，再研究轻量信息价值估计及控制策略；只有开发证据支持时才增加训练复杂度。这一安排有助于把算法效果与系统调试问题分开。

\textbf{人员与组织可行性。}申请人负责科学问题、核心方法和实验质量；拟按状态压缩、恢复校验、评测工程三个任务组组织研究生协作，各组共享版本化任务清单和错误分析记录。本项目实际参与人员与投入安排为\textbf{【待补 F-KS-02：研究生团队名单、分工与投入期限】}。其它申请书中已确认的团队名单不自动转用为本项目成员，核心实验先按现有条件和明确负责人可组织的范围推进。

\textbf{进度可行性。}一年期草案以公开环境的机制实验和轻量原型为主体，企业场景作为授权后的扩展。跨环境测试使用已有任务体系，不新增大规模基座训练或自建综合业务平台。各阶段以可运行基线、错误可定位和结果可复现验收；若方法假说未得到支持，仍可形成边界分析与比较证据。

\subsubsection{研究风险与应对措施}
\textbf{状态提取错误导致依赖保护无效。}早期信号是人工核对发现未完成目标与已完成动作混淆，或依赖保留方案在相同长度下不优于近因规则。预案是收缩到有明确接口前提的任务类型，保留关键原文和确定性规则，先完成可观测条件下的机制实验；对自然语言开放依赖只报告探索性结果。

\textbf{版本信息不完备使恢复不能可靠判断。}早期信号是新旧来源无法区分、元数据冲突或检索命中后仍频繁触发动作前提违例。预案是对当前状态可查询的任务使用主动重查，单独分析不可查询的情境；不能恢复确定证据时结束本次任务并计为失败。完整性哈希与来源版本分开解释，避免用完整性检查代替有效性判断。

\textbf{额外控制开销抵消收益。}早期信号是窗口峰值下降，但累计调用或端到端时延上升。预案是限制压缩频率和恢复循环，比较更简单的规则与轻量压缩器，并将短任务保留历史作为合理选择。若净收益只在高噪声长任务存在，就明确报告该范围，不将结果推广到所有任务。

\textbf{近期研究重合与复现资源不足。}若强组合基线已达到同等表现，将按证据撤回独有机制或性能优越的表述；若原生训练型方法超出现有资源，采用同骨干机制复现并标注差异，不把未复现方法作为弱基线。核心成果可转向元数据可观测性下的可靠性与成本边界，为后续研究保留可复核依据。

\textbf{企业合作条件影响外部验证。}若计划中的数据授权、接口或协作人员未落实，先完成公开可复位环境、构造压力测试及迁移要求说明。企业数据不作为主要机制结论的唯一来源，且不把公开构造任务声称为真实业务数据。授权落实后再按相同评价口径补充外部验证，不反向调整主测试结论。
